\defmodule{BASISLR}

This module offers basic tools to manipulate scalars, vectors, and matrices
that are used to represent lattice bases, using the {\tt LONGREAL} type.
It is used by the modules {\tt LATBASIS}, {\tt REDBAS}, etc.,
when this representation type is used.
See that module for more details.

\bigskip\hrule
\code

DEFINITION MODULE BASISLR;

FROM SUPINT   IMPORT SuperInteger;
FROM CONFIG   IMPORT BMaxDim, BDim;

TYPE
   BScal   = LONGREAL;
   BVect   = ARRAY BDim OF BScal;
   BMatr   = ARRAY BDim OF BVect;
   VectLI  = ARRAY BDim OF LONGINT;

VAR
   BZero, BOne, BMinusOne : BScal;    (* 0, 1, and -1.                   *)

PROCEDURE BCreateScal (VAR S : BScal);

PROCEDURE BDeleteScal (VAR S : BScal);

PROCEDURE BCreateVect (VAR V : BVect; d1, d2 : BDim);

PROCEDURE BDeleteVect (VAR V : BVect; d1, d2 : BDim);

PROCEDURE BCreateMatr (VAR M : BMatr; d1, d2 : BDim);

PROCEDURE BDeleteMatr (VAR M : BMatr; d1, d2 : BDim);
\endcode
  \tab These 6 procedures do nothing.
  There are there only for compatibility with the corresponding
  procedures in {\tt BASISSI}.
  \endtab
\code

PROCEDURE BReadScal (VAR S : BScal);

PROCEDURE BWriteScal (S : BScal;   nb : CARDINAL);
\endcode
  \tab  Writes content of {\tt S} in current output file.
        Uses a total of at least {\tt nb} positions.
  \endtab
\code

PROCEDURE BWriteVect (VAR V : BVect;  n : BDim;  nb : CARDINAL);
\endcode
  \tab  Writes content of {\tt V} for indices {\tt [1..n]}
        in current output file.
        Uses a total of at least {\tt nb} positions for each number.
  \endtab
\code

PROCEDURE BWriteMatr (VAR M : BMatr;  n : BDim; nb : CARDINAL);
\endcode
  \tab  Writes content of {\tt M} for indices {\tt [1..n][1..n]}
        in current output file.
        Uses a total of at least {\tt nb} positions for each number.
  \endtab
\code

PROCEDURE BProdScalVect (s : BScal;  VAR V : BVect;  d : BDim);
\endcode
  \tab Multiplies vector {\tt V} by scalar {\tt s}. \endtab
\code

PROCEDURE BProdScal (VAR U, V : BVect;  n : BDim;  VAR S : BScal);
\endcode
   \tab Returns in {\tt S} the scalar product of {\tt U} by {\tt V},
    using components from 1 to {\tt n}. \endtab
\code

PROCEDURE BProdScalLR (VAR U, V : BVect;  n : BDim) : LONGREAL;
\endcode
   \tab Returns the scalar product of {\tt U} by {\tt V},
    using components from 1 to {\tt n}. \endtab
\code

PROCEDURE BProdScalSI (VAR U, V : BVect; n : BDim; Res : SuperInteger);
\endcode
  \tab  Returns in {\tt S} the scalar product of {\tt U} by {\tt V},
    using components from 1 to {\tt n}.  {\tt U} and {\tt V} are
    rounded and converted into {\tt SuperInteger} representation before
    the computation of the scalar product. \endtab
\code

PROCEDURE BL1NormLR (VAR V : BVect;  n : BDim) : LONGREAL;
\endcode
   \tab Returns the $L_1$ norm of {\tt V},
    using components from 1 to {\tt n}. \endtab
\code

PROCEDURE BL2NormLR (VAR V : BVect;  n : BDim) : LONGREAL;
\endcode
   \tab Returns the $L_2$ norm of {\tt V},
    using components from 1 to {\tt n}. \endtab
\code

PROCEDURE BL1NormSI (VAR V : BVect;  n : BDim;  S1 : SuperInteger);
\endcode
   \tab Returns in {\tt S1} the $L_1$ norm of {\tt V},
    using components from 1 to {\tt n}. \endtab
\code

PROCEDURE BSupNormLR (VAR V : BVect;  n : BDim) : LONGREAL;
\endcode
   \tab Returns the sup norm of {\tt V},
    using components from 1 to {\tt n}. \endtab
\code

PROCEDURE BSupNormSI (VAR V : BVect;  n : BDim;  S1 : SuperInteger);
\endcode
   \tab Returns in {\tt S1} the sup norm of {\tt V},
    using components from 1 to {\tt n}. \endtab
\code

PROCEDURE BModifVectLR (VAR U, V : BVect;  c : LONGREAL;  n : BDim);
\endcode
  \tab  Adds {\tt c} times the vector {\tt V} to the vector {\tt U}.
   Assumes that both vectors have dimension {\tt n}.\endtab
\code

PROCEDURE BModifVectLI (VAR U, V : BVect;  c : LONGINT;  n : BDim);
\endcode
  \tab  Adds {\tt c} times the vector {\tt V} to the vector {\tt U}.
   Assumes that both vectors have dimension {\tt n}.\endtab
\code

PROCEDURE BModifVect (VAR U, V : BVect;  c : BScal;  n : BDim);
\endcode
  \tab  Adds {\tt c} times the vector {\tt V} to the vector {\tt U}.
   Assumes that both vectors have dimension {\tt n}.\endtab
\code

PROCEDURE BChangeSign (VAR S : BScal);
\endcode
  \tab  Changes the sign of S. \endtab
\code

PROCEDURE BChangeSignVect (VAR V : BVect;  n : BDim);
\endcode
  \tab  Changes the sign of elements {\tt 1} to {\tt n} of vector {\tt V}.
  \endtab
\code

PROCEDURE BAffect (S : BScal;  VAR T : BScal);
\endcode
  \tab  Gives to {\tt T} the value of {\tt S}.\endtab
\code

PROCEDURE BExchange (VAR S, T : BScal);

PROCEDURE BEqual (S, T : BScal): BOOLEAN;
\endcode
  \tab  Returns {\tt TRUE} if {\tt S} is equal to {\tt T},
   {\tt FALSE} otherwise.
  \endtab
\code

PROCEDURE BSmaller (S, T : BScal): BOOLEAN;
\endcode
  \tab  Returns {\tt TRUE} if {\tt S} is smaller than {\tt T},
  {\tt FALSE} otherwise.
  \endtab
\code

PROCEDURE BGreater (S, T : BScal): BOOLEAN;
\endcode
  \tab  Returns {\tt TRUE} if {\tt S} is greater than {\tt T},
  {\tt FALSE} otherwise.
  \endtab
\code

PROCEDURE BAbsSizeOrd (S, T : BScal): INTEGER;
\endcode
  \tab  Returns -1 (resp. 0, +1) according as {\tt S} $>$ (resp. $=, <$) {\tt T}. 
  \endtab
\code

PROCEDURE BAbs (S : BScal;  VAR T : BScal);
\endcode
  \tab  Returns in {\tt T} the absolute value of {\tt S}. \endtab
\code

PROCEDURE BSetNeg (VAR S : BScal);
\endcode
  \tab  Changes the sign of {\tt S} if it is not negative.
  \endtab
\code

PROCEDURE BInc (VAR S : BScal; i : LONGINT);
\endcode
  \tab  Adds  {\tt i}  to  {\tt S}.
  \endtab
\code

PROCEDURE BAdd (S, T : BScal;  VAR Res : BScal);
\endcode
  \tab  Returns in {\tt Res} the sum of {\tt S} and {\tt T}.  \endtab
\code

PROCEDURE BMultiply (S, T : BScal;   VAR Res : BScal);
\endcode
  \tab  Returns in {\tt Res} the product of {\tt S} by {\tt T}. \endtab
\code

PROCEDURE BDivideRound (S, T : BScal; VAR Res : BScal);
\endcode
  \tab  Returns in {\tt Res} the result of the division of {\tt S}
        by {\tt T}.  The result is rounded.  For results just in the
        middle of two integers, the value returned by the function is
        the integer the nearest of zero.
        Ex.: If $S=5.0$ and $T=10.0$, then the returned value is 0.
            If $S=-5.0$ and $T=10.0$, then the returned value is also 0.
  \endtab
\code

PROCEDURE BCopyVect (VAR U, V : BVect;  n : BDim);
\endcode
  \tab  Copies the vector {\tt U}, of dimension {\tt n}, into vector {\tt V}.
  \endtab
\code

PROCEDURE BCopyMatr (VAR M1, M2 : BMatr;  n : BDim);
\endcode
  \tab  Copies matrix {\tt M1}, of dimension {\tt n}, into matrix {\tt M2}.
  \endtab
\code

PROCEDURE BConvLR (S : BScal): LONGREAL;
\endcode
  \tab  Returns the {\tt LONGREAL} representation of {\tt S}.
  Is there only for compatibility with {\tt BASISSI}. \endtab
\code

PROCEDURE BConvSI (R : BScal;  S : SuperInteger);
\endcode
  \tab  Returns in {\tt S} the {\tt SuperInteger} representation of {\tt R}.
  \endtab
\code

PROCEDURE LIToBScal (i : LONGINT; VAR S : BScal);
\endcode
   \tab Converts i to  {\tt S}\endtab
\code

PROCEDURE BAffectIdentMatr (VAR M : BMatr;  n : BDim);
\endcode
  \tab  Initializes {\tt M[1..n,1..n]} to the identity matrix.\endtab
\code

PROCEDURE BTrMatr (M : BMatr;   VAR Mtr : BMatr;   d : BDim);
\endcode
  \tab  Computes the transposed matrix of {\tt M}, and
        returns the result in {\tt Mtr}.
  \endtab
\code

PROCEDURE BCanonicalVect (VAR V : BVect;  i : BDim);
\endcode
   \tab Produces the vector with all coordinates equal to zero with the
   exception of the $i$-th coordinate which is equal to one.
   \endtab
\code

PROCEDURE BRandomScal (VAR s : BScal; m : BScal);
\endcode
   \tab Produces a random integer in the interval $[0,m)$, with a uniform
   distribution. \endtab
\code

PROCEDURE GCD2vect (VAR V : VectLI;   k, n : BDim) : LONGINT;
\endcode
 \tab  Computes and returns the greatest common divisor of the vectors
  {\tt V[k],\dots,V[n]}.
 \endtab
\code

PROCEDURE BEuclide (A, B : BScal;  VAR C, D, E, F, G : BScal);
\endcode
  \tab For given $A$ and $B$, returns $C$, $D$, $E$, $F$, $G$ such that: 
   \begin {eqnarray*}
     C  A + D  B  &=& G = \mbox{gcd} (A,B) \\
     E  A + F  B  &=& 0,
   \end {eqnarray*}
where {gcd}  is the greatest common denominator.
  \endtab
\code

END BASISLR.
\endcode
